For a trace state \(\kappa\), let \(r(\kappa)\) be the number of random vertices in its current CADMG, and let \(d(c)\) be the number of missingness indicators assigned by its partial context.  Define
\[
\mu(\kappa)=r(\kappa)+|\mathcal R|-d(c).
\tag{1}
\]
This is a nonnegative integer.  Initially \(r(\kappa)\leq n\) and \(|\mathcal R|-d(c)\leq n\), so \(\mu(\kappa)\leq2n\).  Under \(\text{def:tcf-procedure}\), a strict ancestral marginalization, restriction to a proper district child, or licensed fixing reduces the random-vertex set by at least one.  A delayed selection strictly extends \(c\), and hence increases \(d(c)\) by at least one.  No permitted recursive transition increases either summand in (1).  Therefore every recursive operation word has length at most
\[
2n.
\tag{2}
\]

There are at most \(2^n\) observed margins.  For each indicator \(R_Z\), a partial context either leaves it unassigned or assigns one of \(|\mathcal X_{R_Z}|\) values; therefore the number of partial contexts is exactly bounded by
\[
c=\prod_{R_Z\in\mathcal R}(|\mathcal X_{R_Z}|+1).
\tag{3}
\]
Consequently there are at most \(2^n c\) initial traces.  From any state, each of the two strict subset operations has at most \(2^n\) choices, fixing has at most \(n(\ell+1)\) vertex-license choices, selection has at most \(c\) strict extensions, and one further symbol records termination.  Hence the operation alphabet has size at most
\[
L_{\mathrm{op}}=2^{n+1}+n(\ell+1)+c+1.
\tag{4}
\]
For a fixed initial trace and operation word, all components \((\mathcal H,F,c,q,\pi,\Lambda)\) are prescribed by the update rules.  From (2)--(4), the number of literal complete states visited is at most
\[
2^n c\sum_{j=0}^{2n}L_{\mathrm{op}}^j
=B_{\mathrm{state}}.
\tag{5}
\]
This count keeps every fixing-selection interleaving and every ordered provenance list.  By \(\text{lem:colored-adjacent-swap}\), ordinary kernel commutation alone does not identify unequal trace suffixes.  The optional optimization in \(\text{def:tcf-procedure}\) merges branches only after canonical replay and literal equality of the full tuple, so it can only decrease the bound (5).

The response-fiber feasibility and all coordinate-constancy formulas are formed once, independently of the traces.  A doubled constancy formula uses at most \(2\{g+hK_{\mathrm{ub}}(\mathcal G)\}\) response coordinates, and the free observed vector contributes \(|\mathcal X_{\mathcal O}|\) coordinates.  Thus the base input has at most
\[
p=2\{g+hK_{\mathrm{ub}}(\mathcal G)\}+|\mathcal X_{\mathcal O}|
\tag{6}
\]
variables.  Its complete binary encoding length is at most \(s\) by \(\text{def:termination-budget}\).  By the definition of the fixed finite CAD step bound \(\operatorname{CAD}(m,b)\), the single base call therefore terminates within
\[
C_{\mathrm{base}}=\operatorname{CAD}(p,s)
\tag{7}
\]
projection, real-root-isolation, and sign-evaluation steps and output cells.  In particular, the trace-independent partition has at most \(C_{\mathrm{base}}\) cells.

Each visited trace offers at most one candidate expression for each \((a,y)\) and each base cell.  Equations (5) and (7) therefore give at most
\[
N_{\mathrm{check}}
=|\mathcal X_A||\mathcal X_Y|B_{\mathrm{state}}C_{\mathrm{base}}
\tag{8}
\]
whole-cell candidate-equality tests.  After supported denominators are cleared, such a sentence contains one response vector, the scalar \(z\), and \(\mathbf p\).  It consequently has at most \(p+1\) variables; its complete encoding length is at most \(s\) by the definition of that size parameter.  Each test is thus decided within \(\operatorname{CAD}(p+1,s)\) steps.  Summing the trace visits, the one base CAD, and all annotation tests yields
\[
\begin{aligned}
B_{\mathrm{state}}+C_{\mathrm{base}}
+N_{\mathrm{check}}\operatorname{CAD}(p+1,s)
&=B_{\mathrm{tot}}.
\end{aligned}
\tag{9}
\]

Finally, the completed partition and its annotations mention only finitely many integer-polynomial signs.  Under \(\operatorname{Exact}(P_{\mathrm{obs}})\), rational cells are evaluated by exact rational arithmetic, algebraic cells by the supplied square-free polynomials and isolating intervals, and general exact-real cells by the supplied sign oracle.  Hence the unique cell is selected after finitely many sign decisions.  The finite root-isolation output supplies the stated Thom encodings over the real closure of the ordered field generated by \(\mathbf p\).  These specialization operations terminate separately from (9).  Nothing in the argument defines computation on an unencoded population real or bounds (9) by a polynomial.